\subsection{A structural obstruction for every act of a monoid of order at most seven}

\begin{lemma}[The homology of a triangular act cone]
\label{C:act-cone-Tor}
Let $M$ be a finite monoid, let $S$ be a nonempty finite right
$M$-set, and let $\mathcal C(M,S)$ have two objects, endomorphism
monoids $M$ and $\{1\}$, and $\Hom(0,1)=S$, with composition given
by the action. Put $R=\mathbb QM$ and
$L=\ker(\mathbb QS\xrightarrow{\epsilon}\mathbb Q)$.
For every $n\geq1$ there is a natural isomorphism
\[
 \widetilde H_n(B\mathcal C(M,S);\mathbb Q)
 \cong\operatorname{Tor}^{R}_{n-1}(L,\mathbb Q).
\]
The same conclusion, with $L$ replaced by its coinvariants for the
other vertex monoid, holds for a triangular biset category with
endomorphism monoids $M,N$, provided the trivial left
$\mathbb QN^{\mathrm{op}}$-module is projective.
\end{lemma}
\begin{proof}
The mixed simplices in the nerve give the homotopy pushout
\[
 BM\ \longleftarrow\ E_S\ \longrightarrow\ *,
 \qquad E_S=B(S,M,*).
\]
Choose a projective resolution $P_\bullet\to\mathbb Q$ over $R$.
The edge space and the vertex space are computed, naturally with
respect to their map, by
$\mathbb QS\otimes_R P_\bullet$ and
$\mathbb Q\otimes_R P_\bullet$.
Since each $P_q$ is projective, tensoring the exact sequence
$0\to L\to\mathbb QS\to\mathbb Q\to0$ gives a short exact
sequence of complexes. The mapping cone of its last map is
quasi-isomorphic to $(L\otimes_R P_\bullet)[1]$.
The degree-zero augmentation contributes the one unreduced
$\mathbb Q$ in the classifying space and no positive homology.
This proves the displayed formula, including $n=1$.

For a biset, use $R\otimes\mathbb QN^{\mathrm{op}}$ and the
two-sided bar. Projectivity of the trivial module for the second
factor makes its coinvariant functor exact. Thus a projective
resolution for the first factor tensored with that trivial module
is a projective resolution for the product algebra, and taking
the second coinvariants gives the stated formula. In particular
$BN$ is rationally acyclic. Replacing its vertex in the double
mapping cylinder by a point is a rational homology equivalence,
as follows directly from the homotopy-pushout exact sequence.
\end{proof}

\begin{proposition}[No projective-plane act cone over at most seven elements]
\label{C:seven-act-cones}
For every finite monoid $M$ with $|M|\leq7$ and every nonempty
finite right $M$-set $S$, the category $\mathcal C(M,S)$ does not
have classifying space of type $\mathbb CP^2$.
There is no bound on $|S|$ in this assertion. The same exclusion
holds for a triangular biset category with other endomorphism
monoid $N$ whenever the trivial left
$\mathbb QN^{\mathrm{op}}$-module is projective.
\end{proposition}
\begin{proof}
If $H^2(BM;\mathbb Q)=0$, every degree-two class in the
triangular category lifts to cohomology relative to $BM$.
After replacing the rationally acyclic other vertex by a point,
the quotient by $BM$ is a suspension. Relative products therefore
vanish, and so do the products of these degree-two classes in
absolute cohomology. This excludes $\mathbb CP^2$.

For $|M|\leq6$ and $H^2(BM;\mathbb Q)\ne0$, the proof of
Lemma~\ref{C:small-triangular-vertices} supplies either a projective
resolution of the trivial rational module of length two, or an
exact resolution with modules periodic from degree two and
differentials periodic from degree three. That proof does not
bound the size of the act. We now verify the same sufficient
obstruction for every seven-element target, with one additional
syzygy pattern.

The classification of the $43$ seven-element monoids with
nonzero rational $H^2$ is the complete translation-table
enumeration in Lemma~\ref{C:triangular-eleven}. In particular it
includes reducible monoids as well as the irreducible cores; no
homotopy equivalence of monoids is silently used to transfer an
arbitrary act. Number these tables $0,\ldots,42$ in the
deterministic order of \texttt{a\_CP2\_11\_triangular.targets()}.
The exact resolution computation gives the following partition:
\[
\begin{array}{c|r|l}
 \text{property}&\text{count}&\text{table indices}\\\hline
 \operatorname{pd}_{\mathbb QM}\mathbb Q\leq2&27&
 3,4,5,6,8,9,10,11,14,15,16,17,18,19,22,24,\\
 &&27,28,29,32,33,34,35,36,37,41,42\\
 \text{two-periodic resolution tail}&11&
 1,2,7,13,21,23,26,31,38,39,40\\
 \Omega^3\mathbb Q\cong(\Omega^2\mathbb Q)^{\oplus2}&4&0,12,25,30\\
 \text{central }C_2\text{ splitting}&1&20
\end{array}
\]
Here is what is checked, so that these are all-degree certificates.
The first $27$ resolutions are integral resolutions by summands
$\mathbb ZM e$ and $\mathbb ZM(e-f)$, with
$e^2=e$, $f^2=f$, $ef=fe=f$. Every differential is equivariant,
every image is saturated over $\mathbb Z$, and the successive
rank sums give exactness. The final degree-two differential is
injective. The separate program
\texttt{b\_verify\_resolutions.py} checks these facts independently
of the resolution search.

For the $11$ periodic cases, the projectives in every degree from
two onward are the same explicit module. The matrices satisfy
$d_3=d_5$, and $d_4,d_5$ form an exact two-periodic complex:
both products are zero, both images are saturated, and their ranks
sum to the rank of the common projective. Thus the displayed
finite initial segment continues to an exact resolution with
$d_{q+2}=d_q$ for every $q\geq3$.
For every right module $L$ it follows that
\[
 \operatorname{Tor}_{q+2}^{\mathbb QM}(L,\mathbb Q)
 \cong\operatorname{Tor}_{q}^{\mathbb QM}(L,\mathbb Q)
 \qquad(q\geq3).
\]

For each of the four replication cases, the program
\texttt{a\_seven\_cone\_syzygies.py} separately verifies the
projective initial segment through $P_2$, including saturation
and exactness at $P_0,P_1$. It computes the actual kernels
$\Omega^2\mathbb Q$ and $\Omega^3\mathbb Q$, of dimensions
four and eight, with the action of every monoid element.
It then supplies an $8$ by $8$ rational intertwining matrix
from two copies of the former module to the latter. The
determinants, in index order $0,12,25,30$, are respectively
$-2,-6,-2,8$. Every intertwining equation is verified exactly;
all matrices and module actions are in the JSON certificate.
Dimension shifting therefore gives, for every right module $L$,
\[
 \operatorname{Tor}_{q+1}^{\mathbb QM}(L,\mathbb Q)
 \cong
 \operatorname{Tor}_{q}^{\mathbb QM}(L,\mathbb Q)^{\oplus2}
 \qquad(q\geq3).
\]

Finally, table $20$ has a central unit $g=6$ of order two which
acts as the identity on every nonunit. Over $\mathbb Q$ the
central idempotents $(1+g)/2$ and $(1-g)/2$ split its algebra
as the product of the six-element Sweeney monoid algebra and
$\mathbb Q$. The first factor has basis its new identity,
the four band elements, and the remaining nonunit; multiplication
is exactly the six-element table. The trivial module is supported
on this first factor. Its Tor groups with an arbitrary module
therefore have the same periodicity just proved for that monoid.

In every case, Lemma~\ref{C:act-cone-Tor} shows that nonzero
fourth rational homology in the triangular category either is
impossible or persists in arbitrarily high degrees: the periodic
cases have $H_{n+2}\cong H_n$ for $n\geq4$, and the replication
cases have $H_{n+1}\cong H_n^{\oplus2}$ for $n\geq4$.
This is incompatible with the rational homology of
$\mathbb CP^2$. Exactness of the second-vertex coinvariants
proves the biset version by the same Tor formulas.

The program \texttt{a\_seven\_cone\_targets.py} produces the
complete resolution records and the partition. The replication matrices
are in \texttt{a\_seven\_cone\_syzygies.json}. The computation
does not treat failure of a resolution search to terminate as
proof of an infinite tail: each used tail has the displayed
exact periodic or module-isomorphism certificate.
\end{proof}
